% Solution to Problem 2 — A universal test vector for Rankin-Selberg integrals (Nelson)
% v8, corrected against First Proof paper (Paul D. Nelson)

\documentclass[11pt]{article}

\usepackage[margin=1in]{geometry}
\usepackage{amsmath, amssymb, amsthm}
\usepackage{mathtools}
\usepackage{xcolor}
\usepackage{hyperref}
\hypersetup{colorlinks=true,
            linkcolor=blue!50!black,
            citecolor=blue!50!black,
            urlcolor=blue!50!black}

\newcommand{\F}{F}
\newcommand{\oo}{\mathfrak{o}}
\newcommand{\pp}{\mathfrak{p}}
\newcommand{\qq}{\mathfrak{q}}
\newcommand{\C}{\mathbb{C}}
\newcommand{\Z}{\mathbb{Z}}
\newcommand{\GL}{\mathrm{GL}}
\newcommand{\Mn}{M_n}
\newcommand{\Se}{\mathcal{S}_{\mathrm e}}
\newcommand{\Sw}{\mathcal{S}}
\newcommand{\W}{\mathcal{W}}
\newcommand{\vol}{\operatorname{vol}}
\newcommand{\diag}{\operatorname{diag}}
\newcommand{\tr}{\operatorname{trace}}
\newcommand{\ttilde}[1]{\widetilde{#1}}
\newcommand{\lRS}{\ell_{\mathrm{RS}}}

\newtheorem{theorem}{Theorem}
\newtheorem{lemma}{Lemma}
\newtheorem{proposition}{Proposition}
\newtheorem{corollary}{Corollary}

\title{Solution to Problem 2: A Universal Test Vector\\
for Rankin--Selberg Integrals}
\author{}
\date{May 27, 2026}

\begin{document}
\maketitle

\section*{Statement and Answer}

Let $\F$ be a non-archimedean local field with ring of integers $\oo$, maximal
ideal $\pp$, residue cardinality $q$, and a nontrivial additive character
$\psi:\F\to\C^\times$ of conductor $\oo$. Write $G_r:=\GL_r(\F)$, let $N_r<G_r$
be the upper-triangular unipotent subgroup, and embed $G_n\hookrightarrow
G_{n+1}$ as the upper-left block. Set $K_r:=\GL_r(\oo)$, with Haar measures
giving $K_r$ and $N_r\cap K_r$ volume one.

\medskip
\noindent\textbf{Theorem.} \emph{Let $\Pi$ be a generic irreducible admissible
representation of $G_{n+1}$, realized in its $\psi^{-1}$-Whittaker model
$\W(\Pi,\psi^{-1})$. Then there exists a single $W\in\W(\Pi,\psi^{-1})$ with the
following property. For every generic irreducible admissible $\pi$ of $G_n$ with
conductor ideal $\qq$, generator $Q$ of $\qq^{-1}$, and
$u_Q:=I_{n+1}+QE_{n,n+1}$, there exists $V\in\W(\pi,\psi)$ such that}
\[
\int_{N_n\backslash G_n} W(\diag(g,1)u_Q)\,V(g)\,|\det g|^{s-\frac12}\,dg
\]
\emph{is finite and nonzero for all $s\in\C$.}

\medskip
\noindent The answer is \textbf{YES}. The single $W$ may be taken to be an
explicit, $\Pi$-independent-of-$\pi$ Kirillov-model vector, and for each $\pi$
the role of $V$ is played by the \emph{normalized newvector}. The integral is in
fact \emph{constant in $s$} and equal to an explicit nonzero scalar; this is why
finiteness and nonvanishing hold for all $s$, with no analytic continuation or
pole-cancellation required.

\medskip
\noindent\textbf{A failure mode to avoid.} One cannot choose $V$ as a
Schwartz bump in the Kirillov model making the integrand constant on its
support: $V$ has a central character, so it cannot be compactly supported modulo
the center. (Already for $n=1$ the relevant integral is a normalized Gauss sum,
whose integrand is genuinely non-constant.) The correct argument runs through the
Godement--Jacquet functional equation, not through test-vector flexibility.

\section*{Choice of $W$ and $V$}

For $x\in\F$ and $y\in\F^\times$ write
\[
d_y:=\diag(1,\dots,1,y)\in G_n\hookrightarrow G_{n+1},\qquad
u_x:=I_{n+1}+xE_{n,n+1}\in N_{n+1}.
\]

\noindent\textbf{The vector $W$.} Define $W_0$ on $G_n$ by
\begin{equation}\label{eq:W0}
W_0(g):=\int_{N_n}\mathbf 1_{K_n}(xg)\,\psi(x)\,dx,
\end{equation}
a $\psi^{-1}$-Whittaker function on $G_n$ which extends, via the theory of the
Kirillov model \cite{Bernstein}, to an element of $\W(\Pi,\psi^{-1})$ on
$G_{n+1}$. We take $W:=u_Q W_0$ (equivalently we evaluate $W_0$ at the shifted
argument); $W_0$ depends only on $\Pi$, not on $\pi$.

\noindent\textbf{The vector $V$.} Given $\pi$ with conductor $\qq$, let
$V\in\W(\pi,\psi)$ be the normalized newvector: the unique $K_1(\qq)$-invariant
vector with $V(1)=1$, where
$K_1(\qq)=\{g\in K_n:\text{last row}\equiv(0,\dots,0,1)\bmod\qq\}$
\cite{JPSS81,Matringe}.

We choose a generator $Q$ of $\qq^{-1}$, so $|Q|=[\oo:\qq]$. Recall the
conductor is characterized by the $\varepsilon$-factor:
\begin{equation}\label{eq:eps}
\varepsilon(\tfrac12+s,\pi,\psi)=|Q|^{-s}\,\varepsilon(\tfrac12,\pi,\psi).
\end{equation}

\section*{The Godement--Jacquet engine}

\begin{lemma}[Godement--Jacquet functional equation \cite{GJ72}]\label{lem:GJ}
Let $f$ be a matrix coefficient of $\pi$ and $\phi\in\Sw(\Mn(\F))$. The zeta
integral
\[
Z(\phi,f,s):=\int_{G_n}\phi(g)f(g)\,|\det g|^{\frac{n-1}{2}+s}\,dg
\]
converges for $\Re(s)$ large, continues meromorphically with
$Z(\phi,f,s)/L(s,\pi)$ holomorphic, and satisfies
\[
\gamma(s,\pi,\psi)\,Z(\phi,f,s)=Z(\widehat\phi,f^\vee,1-s),\qquad
\gamma(s,\pi,\psi)=\varepsilon(s,\pi,\psi)\frac{L(1-s,\ttilde\pi)}{L(s,\pi)},
\]
where $f^\vee(g)=f(g^{-1})$ and $\widehat\phi$ is the $\psi$-self-dual Fourier
transform on $\Mn(\F)$. The identities persist when $f$ is taken to be a
Whittaker function of $\pi$.
\end{lemma}

Let $\Se(\F^\times)$ be the Schwartz--Bruhat functions $\beta$ with
$\beta(xy)=\beta(x)$ whenever $|y|=1$ (i.e.\ $\beta$ depends only on $|x|$),
with Mellin inversion $\beta(y)=\int_{(\sigma)}\tilde\beta(s)|y|^s\,ds$. For such
$\beta$ define $\beta^\sharp$ by
$\beta^\sharp(y)=\int_{(\sigma)}\tilde\beta(s)|y|^{-s}\,ds/\gamma(\tfrac12+s,\pi,\psi)$.

\begin{lemma}\label{lem:beta}
Define $\beta$ by $\tilde\beta(s)=\varepsilon(\tfrac12+s,\pi,\psi)/L(\tfrac12+s,\pi)$.
Then:
\begin{enumerate}\setlength\itemsep{2pt}
\item $\beta$ is supported on $\{|Q|\le|y|\le|Q|q^n\}$ and equals
$\varepsilon(\tfrac12,\pi,\psi)$ on $\{|y|=|Q|\}$;
\item $\beta^\sharp$ has Mellin transform $\widetilde{\beta^\sharp}(s)=1/L(\tfrac12+s,\ttilde\pi)$,
is supported on $\{1\le|y|\le q^n\}$, and equals $1$ on $\oo^\times=\{|y|=1\}$.
\end{enumerate}
\end{lemma}

\begin{proof}
By definition of $\gamma$, $\widetilde{\beta^\sharp}(s)=\tilde\beta(s)/\gamma(\tfrac12+s,\pi,\psi)
=L(\tfrac12-s,\ttilde\pi)^{-1}\cdot(\cdots)$ simplifies to $1/L(\tfrac12+s,\ttilde\pi)$.
The inverse $L$-values are monic polynomials in $q^{-s}$ of degree $\le n$, so by
Mellin inversion and \eqref{eq:eps}, $\beta$ and $\beta^\sharp$ have the stated
supports and boundary values.
\end{proof}

\begin{lemma}\label{lem:transfer}
For $\pi$ unitary generic, with $f$ a Whittaker function and
$\phi\in\Sw(\Mn(\F))$,
\[
\int_{G_n}\phi(g)f(g)\beta(\det g)|\det g|^{n/2}\,dg
=\int_{G_n}\widehat\phi(g)f^\vee(g)\beta^\sharp(\det g)|\det g|^{n/2}\,dg.
\]
\end{lemma}

\begin{proof}
Insert the Mellin expansion of $\beta$ at $\sigma=0$; the resulting double
integral converges absolutely, and is recognized as
$\int_{(0)}\tilde\beta(s)Z(\phi,f,\tfrac12+s)\,ds$. Apply the functional equation
(Lemma~\ref{lem:GJ}) and recognize the result as the right-hand side.
\end{proof}

\section*{The local computation}

Set $t_Q:=d_Q^{-1}u_Q=u_1 d_Q^{-1}$.

\begin{lemma}\label{lem:relations}
There exist $\beta\in\Se(\F^\times)$ and $\phi\in\Sw(\Mn(\F))$ with, for all
$g\in G_n$,
\begin{align}
\int_{N_n}\beta(\det xg)\,\phi(xg)\,\psi(x)\,dx
&=\varepsilon(\tfrac12,\pi,\psi)\,W_0(g\,t_Q),\label{eq:rel1}\\
\beta^\sharp(\det g)\,\widehat\phi(g)&=|Q|^n\,\mathbf 1_{K_1(\qq)}(g).\label{eq:rel2}
\end{align}
\end{lemma}

\begin{proof}
Take $\beta$ as in Lemma~\ref{lem:beta} and
$\phi(x):=\psi(-x_{nn})\phi_0(xd_Q^{-1})$ with $\phi_0=\mathbf 1_{\Mn(\oo)}$. For
\eqref{eq:rel1}, use $W_0(gt_Q)=\psi(-g_{nn})W_0(gd_Q^{-1})$ and the identities
$\phi(xg)=\psi(-g_{nn})\phi_0(xgd_Q^{-1})$,
$\beta(\det g)\phi_0(gd_Q^{-1})=\varepsilon(\tfrac12,\pi,\psi)\mathbf 1_{K_n}(gd_Q^{-1})$
(the latter from Lemma~\ref{lem:beta}(1) and
$\mathbf 1_{K_n}(g)=\mathbf 1_{\oo^\times}(\det g)\phi_0(g)$), then integrate
against $\psi(x)\,dx$ recalling \eqref{eq:W0}. For \eqref{eq:rel2}, the Fourier
computation $\widehat\phi(x)=|Q|^n\widehat{\phi_0}(d_Q(x-E_{nn}))$ with
$\widehat{\phi_0}=\phi_0$, together with Lemma~\ref{lem:beta}(2), gives
$\beta^\sharp(\det g)\widehat\phi(g)=|Q|^n\mathbf 1_{K_1(\qq)}(g)$, since
$K_1(\qq)=\{g:d_Q(g-E_{nn})\in\Mn(\oo),\ \det g\in\oo^\times\}$.
\end{proof}

\section*{Conclusion}

For $W\in\W(\Pi,\psi^{-1})$, $V\in\W(\pi,\psi)$, write the Rankin--Selberg
integral
\[
\lRS(s,W,V):=\int_{N_n\backslash G_n}W(\diag(g,1))\,V(g)\,|\det g|^{s-\frac12}\,dg.
\]

\begin{proposition}\label{prop:final}
Let $W_0$ be as in \eqref{eq:W0} and $V$ the normalized newvector of $\pi$. Then
for all $s\in\C$,
\[
\lRS(s,u_Q W_0,\,d_Q V)=c\,|Q|^{-n/2},\qquad
c:=\varepsilon(\tfrac12,\pi,\psi)^{-1}|Q|^n\vol(K_1(\qq))\asymp 1.
\]
\end{proposition}

\begin{proof}
A change of variables gives the homogeneity
$\lRS(s,u_Q W_0,d_Q V)=|Q|^{-(s-\frac12)}\lRS(s,t_Q W_0,V)$. Since $W_0$ is
supported on $\det^{-1}(\oo^\times)$, the translate $t_Q W_0$ is supported on
$\det^{-1}(Q\oo^\times)$, so $\lRS(s,t_Q W_0,V)$ is a constant multiple of
$|Q|^s$; it therefore suffices to evaluate at $s=\tfrac{n+1}{2}$. Unfolding and
applying Lemmas~\ref{lem:relations} (i.e.\ \eqref{eq:rel1}), \ref{lem:transfer},
and \eqref{eq:rel2} in turn, with $f(g)=V(g)$,
\begin{align*}
\varepsilon(\tfrac12,\pi,\psi)\,\lRS(\tfrac{n+1}{2},t_Q W_0,V)
&=\varepsilon(\tfrac12,\pi,\psi)\!\int_{N_n\backslash G_n}\!W_0(gt_Q)V(g)|\det g|^{n/2}dg\\
&=\int_{G_n}\phi(g)f(g)\beta(\det g)|\det g|^{n/2}dg\\
&=\int_{G_n}\widehat\phi(g)f^\vee(g)\beta^\sharp(\det g)|\det g|^{n/2}dg\\
&=|Q|^n\!\int_{K_1(\qq)}\!V(g^{-1})|\det g|^{n/2}dg
=|Q|^n\vol(K_1(\qq)),
\end{align*}
using $K_1(\qq)$-invariance, $V(1)=1$, and $|\det g|=1$ on $K_1(\qq)$. Hence
$\lRS(\tfrac{n+1}{2},t_Q W_0,V)=c$, so $\lRS(s,t_Q W_0,V)=c\,|Q|^{s-\frac{n+1}{2}}$
and $\lRS(s,u_Q W_0,d_Q V)=c\,|Q|^{-n/2}$.
\end{proof}

Since $c\,|Q|^{-n/2}$ is a nonzero scalar independent of $s$, the integral is
finite and nonzero for all $s\in\C$. The vector $W:=u_Q W_0$ (built from the
single, $\pi$-independent $W_0$) and the $\pi$-newvector $V$ (rescaled by $d_Q$)
exhibit the required test vectors. This proves the Theorem. \qed

\section*{Remarks}

The mechanism is transparent: $u_Q W_0$ is supported on a single $\det$-coset, so
the zeta integral collapses to a \emph{constant} in $s$ rather than an
$L$-factor times a unit. The Godement--Jacquet functional equation, packaged
through the $\gamma$-factor transform $\beta\mapsto\beta^\sharp$
(Lemma~\ref{lem:beta}), converts the $\psi$-twisted unipotent shift into an
integral over a translate of $K_1(\qq)$, where newvector theory yields
$\vol(K_1(\qq))$. The $n=1$ case recovers the classical normalized Gauss sum and
has appeared in subconvexity and moment applications \cite{Sarnak,MV}; the
general $n$ statement, with a single unipotent translate rather than an average,
is the present lemma (joint work in progress with S.~Jana). A version with an
average over many translates appears in \cite{BKL}.

\begin{thebibliography}{9}
\bibitem{Bernstein} J.~N.~Bernstein.
  \emph{$P$-invariant distributions on $\GL(N)$ and the classification of
  unitary representations of $\GL(N)$ (non-archimedean case)}. In
  \emph{Lie group representations II}, LNM \textbf{1041}, Springer, 1984,
  50--102.
\bibitem{BKL} A.~R.~Booker, M.~Krishnamurthy, M.~Lee.
  \emph{Test vectors for Rankin--Selberg $L$-functions}.
  J.\ Number Theory \textbf{209} (2020), 37--48.
\bibitem{GJ72} R.~Godement, H.~Jacquet.
  \emph{Zeta functions of simple algebras}. LNM \textbf{260}, Springer, 1972.
\bibitem{JPSS81} H.~Jacquet, I.~I.~Piatetski-Shapiro, J.~Shalika.
  \emph{Conducteur des repr\'esentations du groupe lin\'eaire}.
  Math.\ Ann.\ \textbf{256} (1981), 199--214.
\bibitem{Matringe} N.~Matringe.
  \emph{Essential Whittaker functions for $\GL(n)$}.
  Doc.\ Math.\ \textbf{18} (2013), 1191--1214.
\bibitem{MV} P.~Michel, A.~Venkatesh.
  \emph{The subconvexity problem for $\GL_2$}.
  Publ.\ Math.\ IH\'ES \textbf{111} (2010), 171--271.
\bibitem{Sarnak} P.~Sarnak.
  \emph{Fourth moments of Gr\"ossencharakteren zeta functions}.
  Comm.\ Pure Appl.\ Math.\ \textbf{38} (1985), 167--178.
\end{thebibliography}

\end{document}